% ============================================================
%  Programación Científica 2026-1 — UNAL
%  Hoja de pistas A4: Quiz de Autodiff
% ============================================================
\documentclass[11pt]{article}

\usepackage[a4paper,top=1.2cm,bottom=1.2cm,left=1.3cm,right=1.3cm]{geometry}
\usepackage[T1]{fontenc}
\usepackage[utf8]{inputenc}
\usepackage{helvet}
\renewcommand{\familydefault}{\sfdefault}
\usepackage{amsmath,amssymb}
\usepackage{booktabs}
\usepackage{colortbl}
\usepackage{array}
\usepackage{multicol}
\usepackage{tcolorbox}
\tcbuselibrary{skins}
\usepackage{xcolor}
\pagestyle{empty}

% ---------- paleta ----------
\definecolor{accent}{HTML}{6C63FF}
\definecolor{accentteal}{HTML}{1DB89A}
\definecolor{accentamber}{HTML}{F4A623}
\definecolor{accentcoral}{HTML}{E05A4E}
\definecolor{textgray}{HTML}{6B6B80}

% ---------- cajas ----------
\newtcolorbox{graybox}[1]{enhanced,title=#1,
  colback=accent!6,colframe=accent,coltitle=white,
  fonttitle=\bfseries\small,boxrule=0.5pt,
  attach boxed title to top left={yshift=-2mm,xshift=4mm},
  boxed title style={colback=accent,boxrule=0pt},
  left=4pt,right=4pt,top=6pt,bottom=4pt}

\newtcolorbox{tealbox}[1]{enhanced,title=#1,
  colback=accentteal!7,colframe=accentteal,coltitle=white,
  fonttitle=\bfseries\small,boxrule=0.5pt,
  attach boxed title to top left={yshift=-2mm,xshift=4mm},
  boxed title style={colback=accentteal,boxrule=0pt},
  left=4pt,right=4pt,top=6pt,bottom=4pt}

\newtcolorbox{amberbox}[1]{enhanced,title=#1,
  colback=accentamber!8,colframe=accentamber,coltitle=white,
  fonttitle=\bfseries\small,boxrule=0.5pt,
  attach boxed title to top left={yshift=-2mm,xshift=4mm},
  boxed title style={colback=accentamber,boxrule=0pt},
  left=4pt,right=4pt,top=6pt,bottom=4pt}

\newtcolorbox{coralbox}[1]{enhanced,title=#1,
  colback=accentcoral!7,colframe=accentcoral,coltitle=white,
  fonttitle=\bfseries\small,boxrule=0.5pt,
  attach boxed title to top left={yshift=-2mm,xshift=4mm},
  boxed title style={colback=accentcoral,boxrule=0pt},
  left=4pt,right=4pt,top=6pt,bottom=4pt}

\begin{document}

% ============================================================
%  ENCABEZADO
% ============================================================
\begin{center}
  {\Large\bfseries Hoja de pistas -- Autodiff}\\[2pt]
  {\small\textcolor{textgray}{Programación Científica 2026-1 · UNAL}}
\end{center}
\vspace{2pt}
\hrule
\vspace{6pt}

% ============================================================
%  FILA 1: Grafo + Derivadas elementales
% ============================================================
\begin{minipage}[t]{0.36\textwidth}
\begin{amberbox}{Paso 1 — El grafo computacional}
\small
\textbf{Cómo leer el código:}
\begin{itemize}\setlength\itemsep{1pt}
  \item Cada asignación es un \textbf{nodo}.
  \item Arista $u \to v$: $v$ depende de $u$.
  \item Nodos hoja (sin padres): entradas $v_1=x$, $v_2=y$.
  \item Nodo raíz (sin hijos): salida $f$.
\end{itemize}
\vspace{3pt}
\textcolor{accentcoral}{\textbf{Atención — nodo con varios hijos:}}\\
Si \texttt{x} aparece en varias líneas del código, $v_1$ tiene \textbf{varios hijos} en el grafo.\\
El gradiente debe \textbf{sumar} la contribución de cada camino.
\end{amberbox}
\end{minipage}
\hfill
\begin{minipage}[t]{0.60\textwidth}
		\vspace{-5.5cm}
\begin{tealbox}{Derivadas elementales (se usan igual en forward y reverse)}
\small
\renewcommand{\arraystretch}{1.18}
\begin{tabular}{lll@{\hspace{6mm}}lll}
Operación & $\partial v/\partial u$ & & Operación & $\partial v/\partial u$ \\
\midrule
$v = u + w$    & $1$ (para $u$ y $w$)  & & $v = \sin u$  & $\cos u$ \\
$v = u - w$    & $1$ (para $u$), $-1$ (para $w$)  & & $v = \cos u$  & $-\sin u$ \\
$v = u\cdot w$ & $w$ (para $u$), $u$ (para $w$)   & & $v = e^u$     & $e^u$ \\
$v = u^n$      & $n\,u^{n-1}$         & & $v = \ln u$   & $1/u$ \\
$v = -u$       & $-1$                 & & $v = c\,u$    & $c$ \\
\end{tabular}
\end{tealbox}
\end{minipage}

\vspace{8pt}
\hrule
\vspace{6pt}

% ============================================================
%  FILA 2: Forward + Reverse en paralelo
% ============================================================
\begin{minipage}[t]{0.475\textwidth}
\begin{graybox}{Paso 2 — Forward mode (cédula PAR)}
\small
\textbf{Objetivo:} propagar tangentes $\dot{v}_i = \dfrac{\partial v_i}{\partial x_j}$ de entrada a salida.

\vspace{4pt}
\textbf{Inicializar} (una corrida por variable de entrada):
\[
\frac{\partial f}{\partial x}:\;\dot{v}_1=1,\;\dot{v}_2=0
\qquad
\frac{\partial f}{\partial y}:\;\dot{v}_1=0,\;\dot{v}_2=1
\]

\textbf{Propagar} en orden natural (entrada $\to$ salida):
\[
\dot{v}_j = \sum_{k\,\in\,\text{padres}(j)} \frac{\partial v_j}{\partial v_k}\;\dot{v}_k
\]

\vspace{4pt}
\textbf{Leer resultado:} $\dot{v}_\text{salida} = \partial f/\partial x_j$

\vspace{3pt}
\textcolor{textgray}{\small 2 corridas (una por entrada) dan $\nabla f$ completo.}
\end{graybox}
\end{minipage}
\hfill
\begin{minipage}[t]{0.475\textwidth}
\begin{coralbox}{Paso 2 — Reverse mode (cédula IMPAR)}
\small
\textbf{Objetivo:} propagar adjuntos $\bar{v}_i = \dfrac{\partial f}{\partial v_i}$ de salida a entradas.

\vspace{4pt}
\textbf{Paso 2a — Forward:} evaluar y \textbf{guardar} todos los $v_i$.

\vspace{4pt}
\textbf{Paso 2b — Backward:}\\
Inicializar $\bar{v}_\text{salida} = 1$. Recorrer en orden \textbf{inverso}:
\[
\bar{v}_i \;\mathrel{+}=\; \bar{v}_j \cdot \frac{\partial v_j}{\partial v_i}
\qquad \forall\, j \in \text{hijos}(i)
\]
\textcolor{accentcoral}{Sumar sobre \textbf{todos} los hijos de $v_i$.}

\vspace{4pt}
\textbf{Leer resultado al llegar a entradas:}
\[
\bar{v}_1 = \frac{\partial f}{\partial x}, \qquad \bar{v}_2 = \frac{\partial f}{\partial y}
\]
\textcolor{textgray}{\small 1 sola corrida da $\nabla f$ completo.}
\end{coralbox}
\end{minipage}

\vspace{6pt}
\hrule
\vspace{4pt}

% ============================================================
%  FILA 3: recordatorio final
% ============================================================
\begin{center}
\small
\textcolor{textgray}{
\textbf{Forward y reverse producen el mismo gradiente.} \quad
Forward: $n$ corridas para $n$ entradas. \quad
Reverse: 1 corrida da todo $\nabla f$. \quad
Ambos usan las mismas derivadas elementales.
}
\end{center}

\end{document}
